\begin{table}[t]
\centering
\footnotesize
\caption{
Paired effect summary against the strongest aligned baseline. Gains are
reported in percentage points. The effect size $d_z$ is paired Cohen's $d$,
computed as the mean paired gain divided by the standard deviation of the
split-level paired gains. The final row treats the six dataset-level mean
gains as the paired units.
}
\label{tab:paired-effect}
\begin{tabular}{lrrr}
\toprule
Dataset & Mean gain & 95\% CI & $d_z$ \\
\midrule
Amazon-Computers & $+1.87$ & $[-0.07,\,+3.81]$ & $0.69$ \\
Amazon-Photo     & $+0.39$ & $[-1.20,\,+1.98]$ & $0.18$ \\
Chameleon        & $+0.92$ & $[-0.14,\,+1.98]$ & $0.62$ \\
Cornell          & $+2.97$ & $[-4.05,\,+10.00]$ & $0.30$ \\
Texas            & $+1.99$ & $[-2.35,\,+6.33]$ & $0.33$ \\
Wisconsin        & $+0.98$ & $[-3.84,\,+5.81]$ & $0.15$ \\
\midrule
Dataset-level aggregate & $+1.52$ & $[+0.54,\,+2.50]$ & $1.62$ \\
\bottomrule
\end{tabular}
\end{table}

All six dataset-level mean gains are positive. To avoid treating the 60
split-level differences as fully independent observations, our primary
significance statement uses the six dataset-level mean gains as the paired
units. A two-sided one-sample $t$-test over the six dataset-level gains gives
$p=0.0105$, a two-sided Wilcoxon signed-rank test gives $p=0.0313$, and a
two-sided sign test for all six gains being positive gives $p=0.0313$. Thus,
the aggregate paired evidence supports a statistically significant improvement
over the strongest aligned baseline. 



\section{Mechanistic Interventions and White-box Mechanism Checks}
\label{sec:ablation-specialization}

A common weakness of ablation studies is that they are run under a different
split or seed protocol from the main experiment. To avoid this issue, we rerun
all ablations using the same split and seed protocol as the final SRC-v16c
main experiment. Each ablation is evaluated over ten matched repeats, uses the
same train/validation/test policy as the corresponding full-model run, selects
its internal hyperparameters by validation only, and evaluates the test set
only after validation selection. Thus, the ablation study is aligned with the
main experimental protocol rather than being a separate two-run diagnostic.

Because \SRC{} is designed as a diagnostic scaffold, the full model should not
be interpreted as the single best specialist for every dataset. Instead, the
full scaffold exposes which signal families and decision mechanisms are active.
The simplified variants in Table~\ref{tab:ablation-specialization} are
therefore not merely ablations; they are controlled mechanistic interventions
that test whether simple structural changes agree with the dataset mechanisms
measured by the atlas.

\begin{table}[t]
\centering
\scriptsize
\caption{
Aligned mechanistic intervention summary. All variants are run with the same
ten split/seed protocol as the final SRC-v16c experiment. Entries are mean
test accuracy in percent. The bottom rows summarize cross-dataset
generalization. Rank is computed within each dataset among the listed variants;
lower is better. The full model is not always the best dataset-specific
specialist, but it has the best average rank, the best worst-case accuracy,
and is the only variant that remains in the top three on all six datasets. The
simplified variants are used as diagnostic interventions, not as a redefinition
of the main benchmark model.
}
\label{tab:ablation-specialization}
\resizebox{\linewidth}{!}{%
\begin{tabular}{lrrrrrrrl}
\toprule
Dataset
& \textbf{Full}
& Raw
& No high-pass
& No $P^3X$
& No sym
& PCA only
& Ridge only
& Best diagnostic variant \\
\midrule
Amazon-Computers & \textbf{78.71} & 63.46 & 79.73 & 76.32 & 73.54 & 76.83 & 78.58 & No high-pass (79.73) \\
Amazon-Photo & \textbf{88.76} & 76.88 & 90.27 & 87.16 & 84.78 & 87.15 & 89.12 & No high-pass (90.27) \\
Chameleon & \textbf{72.48} & 44.71 & 72.02 & 71.67 & 72.35 & 70.61 & 71.62 & Full / near-full (72.48) \\
Cornell & \textbf{75.41} & 77.30 & 71.08 & 75.14 & 74.32 & 66.76 & 77.57 & Ridge only (77.57) \\
Texas & \textbf{86.32} & 82.89 & 83.42 & 85.26 & 85.79 & 77.11 & 86.58 & Ridge only (86.58) \\
Wisconsin & \textbf{84.31} & 88.63 & 79.22 & 85.69 & 83.73 & 78.24 & 83.92 & Raw only (88.63) \\
\midrule
Mean accuracy & \textbf{81.00} & 72.31 & 79.29 & 80.21 & 79.09 & 76.12 & 81.23 & --- \\
Worst-case accuracy & \textbf{72.48} & 44.71 & 71.08 & 71.67 & 72.35 & 66.76 & 71.62 & --- \\
Average rank $\downarrow$ & \textbf{2.33} & 5.00 & 3.67 & 3.83 & 4.50 & 6.00 & 2.67 & --- \\
Top-3 count & \textbf{6/6} & 2/6 & 3/6 & 1/6 & 2/6 & 0/6 & 4/6 & --- \\
\bottomrule
\end{tabular}%
}
\end{table}

The aligned interventions show that SRC-v16c should be interpreted as a
general white-box scaffold rather than as a single dataset-specific specialist.
Some simplified variants can outperform or closely match the full scaffold on
individual datasets: removing high-pass blocks improves the two Amazon graphs,
ridge-oriented variants are strong on Cornell and Texas, and raw-only is
strongest on Wisconsin. However, these specialized gains do not transfer
uniformly across datasets. Raw-only collapses on Chameleon, PCA-only is weak
across the board, and no single specialized variant remains reliable on all
six datasets.

The bottom rows of Table~\ref{tab:ablation-specialization} summarize this
generalist-versus-specialist distinction. The full SRC-v16c scaffold has the
best average rank among the listed variants, the highest worst-case accuracy,
and is the only variant that appears in the top three on every dataset. These
cross-dataset statistics are the key reason we keep the full scaffold as the
main method. In addition, the choice is not based on accuracy alone. Although
the Ridge-only variant attains a slightly higher mean accuracy, it is not an
adequate primary model for the present paper's purpose because it removes the
class-subspace branch that is necessary for the downstream atlas analyses,
especially class-subspace complexity, PCA--Ridge complementarity, and the
geometry-versus-boundary interpretation developed in the later sections.
Dataset-specific variants are useful diagnostics, and sometimes useful
specialists, but the full model is the most stable and analytically complete
white-box base model across the benchmark suite.

This is the intended role of the intervention study. Since the model is
white-box, removing or isolating a component is not merely a damage test; it is
a mechanism check. Amazon graphs are low-pass dominated, so high-pass blocks
can be removed for a specialized improvement. Chameleon requires the combined
multi-hop, high-pass, PCA, and ridge structure, so the full model is strongest
there. WebKB graphs are more boundary- or raw-feature dominated, so simplified
specialists can improve particular datasets. These results preview the atlas
analysis below: once the dataset structure is measured, the same measurements
can explain which algorithmic simplifications or emphases are worth trying.

\section{From Node-Level Atlases to Dataset Fingerprints}
\label{sec:mechanistic-atlas}

We call the aggregated atlas of a dataset an \emph{operational dataset
fingerprint}. It is operational because it is computed by a fixed, named
graph-signal dictionary and by fixed white-box decision modules; it is a
dataset fingerprint because applying the same measurement procedure to every
dataset yields comparable signal compositions, decision geometries, class
complexities, and error shifts. The fingerprint is not a causal claim about
the true data-generating process. It is a reproducible measurement of how a
white-box graph classifier uses the dataset.

For dataset $D$, let $\Eval_D$ denote the evaluation node set used for atlas
reporting; in the retrospective analysis below, this is the test-node set.
Let $\mathcal{C}_{\Train}$ denote the set of classes present in the training
split, and write $C_{\Train}=|\mathcal{C}_{\Train}|$. We define the
dataset-level signal means by
\begin{equation}
R_D
=
\frac{1}{|\Eval_D|}
\sum_{i\in\Eval_D}\Rraw(i),
\qquad
L_D
=
\frac{1}{|\Eval_D|}
\sum_{i\in\Eval_D}\Rlow(i),
\qquad
H_D
=
\frac{1}{|\Eval_D|}
\sum_{i\in\Eval_D}\Rhigh(i),
\label{eq:dataset-signal-means}
\end{equation}
the mean class-subspace complexity by
\begin{equation}
C_D
=
\frac{1}{C_{\Train}}
\sum_{c\in\mathcal{C}_{\Train}} r_c,
\label{eq:dataset-complexity}
\end{equation}
and the PCA--Ridge decision fractions by
\begin{equation}
Q^{\mathrm{ridge}}_D
=
\frac{1}{|\Eval_D|}
\sum_{i\in\Eval_D}
\mathbf{1}\!\left\{
\hat y_i^{\mathrm{ridge}}=y_i,\;
\hat y_i^{\mathrm{pca}}\neq y_i
\right\},
\label{eq:ridge-only-fraction}
\end{equation}
\begin{equation}
Q^{\mathrm{hard}}_D
=
\frac{1}{|\Eval_D|}
\sum_{i\in\Eval_D}
\mathbf{1}\!\left\{
\hat y_i^{\mathrm{ridge}}\neq y_i,\;
\hat y_i^{\mathrm{pca}}\neq y_i
\right\}.
\label{eq:hard-fraction}
\end{equation}
Define the correct and wrong subsets of the evaluation nodes by
\begin{equation}
\Eval_D^{\mathrm{correct}}
=
\{\,i\in\Eval_D:\hat y_i=y_i\,\},
\qquad
\Eval_D^{\mathrm{wrong}}
=
\{\,i\in\Eval_D:\hat y_i\neq y_i\,\}.
\label{eq:correct-wrong-sets}
\end{equation}
For the correctness-dependent high-pass shift, let
\begin{equation}
H_D^{\mathrm{correct}}
=
\frac{1}{|\Eval_D^{\mathrm{correct}}|}
\sum_{i\in\Eval_D^{\mathrm{correct}}}\Rhigh(i),
\qquad
H_D^{\mathrm{wrong}}
=
\frac{1}{|\Eval_D^{\mathrm{wrong}}|}
\sum_{i\in\Eval_D^{\mathrm{wrong}}}\Rhigh(i),
\label{eq:correct-wrong-high}
\end{equation}
and define
\begin{equation}
\Delta H_D
=
H_D^{\mathrm{wrong}}
-
H_D^{\mathrm{correct}}.
\label{eq:highpass-shift}
\end{equation}
We then summarize the operational atlas of dataset $D$ by
\begin{equation}
\mathbf{m}(D)
=
\left[
R_D,\,
L_D,\,
H_D,\,
C_D,\,
Q^{\mathrm{ridge}}_D,\,
Q^{\mathrm{hard}}_D,\,
\Delta H_D
\right].
\label{eq:atlas-fingerprint}
\end{equation}

In the present paper, these quantities are reported on test nodes in order to
characterize how the fixed white-box probe behaves on the dataset after
standard supervised evaluation. This post-evaluation setting is intentional:
the goal here is to measure and diagnose dataset behavior under a fixed probe,
not to claim blind pre-test model selection, prospective model selection, or
dataset-independent design rules. Accordingly, the resulting fingerprint is
meant to provide post-evaluation diagnostic guidance only for later model
analysis and dataset-specific modification on the same measured dataset. If
related diagnostic procedures are later used prospectively, then any
correctness-dependent quantities should be computed only on training or
validation nodes, never on test nodes.

\subsection{Graph-signal fingerprints}

\begin{table}[t]
\centering
\small
\caption{Family-size-adjusted node-level graph signal mixture. Values are average test-node signal-family shares in percent. Family composition is computed by first averaging Fisher-weighted block evidence within each signal family and then normalizing across the three families, thereby removing the trivial family-cardinality prior of the raw/low-pass/high-pass partition.}
\label{tab:signal-mix}
\begin{tabular}{lrrrr}
\toprule
Dataset & $n_{test}$ & Raw & Low-pass & High-pass \\
\midrule
Amazon-Computers & 13252 & 15.32 & 69.76 & 14.92 \\
Amazon-Photo & 7250 & 13.35 & 71.97 & 14.68 \\
Chameleon & 456 & 0.31 & 61.32 & 38.37 \\
Cornell & 37 & 28.12 & 34.99 & 36.89 \\
Texas & 38 & 29.47 & 30.59 & 39.94 \\
Wisconsin & 51 & 34.59 & 32.70 & 32.71 \\
\bottomrule
\end{tabular}
\end{table}


\begin{figure}[H]
    \centering
    \includegraphics[width=0.95\linewidth]{fig_signal_simplex.pdf}
    \caption{
    Graph-signal simplex. Each dataset is positioned in a ternary coordinate system
    using its family-size-adjusted node-level raw, low-pass, and high-pass signal shares.
    Marker color encodes high-pass reliance, while callouts report the high-pass and low-pass
    percentages. Amazon graphs remain low-pass dominated after adjustment, whereas
    Chameleon and the WebKB datasets move toward a more mixed raw/high-pass regime,
    showing that the datasets differ not only in accuracy but also in the graph-signal
    mechanisms used by \SRC{}.
    }
    \label{fig:signal-simplex}
\end{figure}
\FloatBarrier

Table~\ref{tab:signal-mix} and Figure~\ref{fig:signal-simplex} report family-size-adjusted
signal-family shares, so the reported composition is not mechanically driven by
the unequal 1:5:3 block counts of the raw/low-pass/high-pass partition.
\dataset{Amazon-Computers} and \dataset{Amazon-Photo} remain low-pass dominated,
but at a corrected level of roughly $70$--$72\%$ rather than the unadjusted
$84$--$85\%$. \dataset{Chameleon} shifts toward a mixed low-pass/high-pass
profile, while \dataset{Cornell}, \dataset{Texas}, and \dataset{Wisconsin}
move into a substantially more raw- and high-pass-balanced regime. Thus, the
classical homophily/heterophily dichotomy can be refined into a measurable
signal composition: a graph may still be low-pass dominated, yet retain a
nontrivial raw or high-pass component after correcting for family cardinality.

\subsection{Dense node-level signal phase portraits}

\begin{figure}[H]
    \centering
    \includegraphics[width=0.98\linewidth]{fig_node_signal_phase_portraits.pdf}
    \caption{Dense node-level signal phase portraits. Each test node is positioned by family-size-adjusted low-pass share (x-axis) and family-size-adjusted high-pass share (y-axis). Blue semi-transparent points show all test nodes and orange semi-transparent points show errors. The plots show where mistakes occur inside the signal space rather than only how many errors occur.}
    \label{fig:signal-phase}
\end{figure}
\FloatBarrier

Figure~\ref{fig:signal-phase} uses all test nodes rather than dataset averages.
The portrait is deliberately not a feature-attribution bar plot: it is a phase
portrait of the family-size-adjusted graph-signal state induced by the explicit dictionary in
Eq.~\eqref{eq:dictionary}. Amazon errors are concentrated away from the
dominant low-pass core. Chameleon occupies a visibly broader mixed
low-pass/high-pass region, consistent with the need to preserve node-neighborhood differences. This is
the key advantage of a named signal dictionary: the space in which a model
succeeds or fails is directly plottable.

\subsection{PCA geometry versus ridge boundary}

\begin{figure}[H]
    \centering
    \includegraphics[width=0.98\linewidth]{fig_decision_margin_phase.pdf}
    \caption{PCA--Ridge decision phase map. Each node is plotted by PCA margin and ridge margin. Blue points are correct predictions and orange points are errors. The axes separate geometry-driven, boundary-driven, easy, and hard cases.}
    \label{fig:decision-phase}
\end{figure}
\FloatBarrier

The PCA residual and ridge boundary are not redundant. Their node-level margins
form a two-dimensional decision phase space (Figure~\ref{fig:decision-phase}).
Nodes in the positive-positive region are easy for both modules;
negative-negative nodes are hard for both; off-axis regions reveal
geometry-only and boundary-only cases. On \dataset{Texas} and
\dataset{Wisconsin}, ridge margins correct a nontrivial fraction of nodes that
the PCA residual alone misses, explaining why the final classifier often
behaves boundary-dominantly on WebKB graphs.

\subsection{Class-subspace geometry complexity}

\begin{figure}[H]
    \centering
    \includegraphics[width=0.90\linewidth]{fig_class_geometry_complexity.pdf}
    \caption{Class-subspace geometry complexity. Each point is one class, positioned by its selected PCA dimension $r_c$ and colored by retained variance. Chameleon requires much larger class subspaces than Amazon and WebKB, indicating that its classes occupy more complex regions of the selected graph-signal space.}
    \label{fig:class-geometry}
\end{figure}
\FloatBarrier

Figure~\ref{fig:class-geometry} gives a class-level view of the PCA component.
Amazon and WebKB classes require relatively small selected dimensions, while
Chameleon approaches the dimension cap for several classes. This supports the
interpretation that Chameleon is not merely a graph-filtering problem: its
class geometry in the selected signal space is intrinsically more complex,
making the ridge boundary an important complement to PCA residuals.

\subsection{Correct-vs-wrong signal shifts}

\begin{figure}[H]
    \centering
    \includegraphics[width=0.84\linewidth]{fig_error_signal_shift.pdf}
    \caption{Error signal shift. The lollipop value is the family-size-adjusted high-pass share of wrong nodes minus the family-size-adjusted high-pass share of correct nodes. Positive values mean errors have more high-pass mass; negative values mean errors lose high-pass mass relative to correct nodes.}
    \label{fig:error-shift}
\end{figure}
\FloatBarrier

Figure~\ref{fig:error-shift} shows that errors are not uniformly distributed
in signal space. On Amazon graphs, errors carry more high-pass mass than
correct nodes, consistent with errors being atypical relative to the dominant
low-pass mechanism. On Chameleon, the shift goes in the opposite direction:
correct nodes rely more on high-pass differences, while errors are relatively
more low-pass. This directly supports the design choice of including
$X-\Prw X$ and $\Prw X-\Prw^2X$ in the dictionary.

\section{From Dataset Fingerprints to Diagnostic Guidance}

\label{sec:atlas-diagnostic-probe}

The mechanistic atlas converts a dataset from a single accuracy number into a
white-box diagnostic object. Instead of asking only which model wins on a
dataset, WG-SRC asks which graph-signal families and decision mechanisms are
associated with success or failure under a fixed white-box probe. This changes
the role of evaluation: the dataset is no longer treated only as a black-box
benchmark, but also as an object whose raw-feature reliance, propagation
reliance, high-pass sensitivity, class geometry, and boundary dependence can be
measured.

The claim supported in this section is deliberately narrower than automatic
algorithm design but stronger than mechanism interpretation alone. We do not
claim that the atlas, by itself, discovers the globally best new model on a
dataset or replaces a full redesign search. Rather, \SRC{} is a dual-purpose
white-box model: it performs prediction and simultaneously produces the
fingerprint $\mathbf{m}(D)$ in Eq.~\eqref{eq:atlas-fingerprint}. That
fingerprint yields mechanism hypotheses---for example, that high-pass blocks
are acting as noise, that raw features must be preserved, or that a
ridge-style boundary is correcting PCA failures. The role of the aligned interventions is to test whether these hypotheses are
supported in the expected direction. In this sense, the atlas is not a blind
redesign oracle, not a prospective model selector, and not a
dataset-independent design rule. Rather, after a standard evaluation of the
fitted white-box scaffold, it provides post-evaluation diagnostic guidance for
later model analysis and dataset-specific modification on that same measured
dataset.


\begin{figure}[t]
\centering
\begin{tikzpicture}[
    node distance=1.2cm and 1.35cm,
    box/.style={
        rectangle,
        rounded corners,
        draw=black!70,
        fill=gray!6,
        thick,
        align=center,
        minimum width=3.0cm,
        minimum height=0.9cm
    },
    keybox/.style={
        rectangle,
        rounded corners,
        draw=blue!70!black,
        fill=blue!6,
        thick,
        align=center,
        minimum width=3.35cm,
        minimum height=0.95cm
    },
    designbox/.style={
        rectangle,
        rounded corners,
        draw=orange!85!black,
        fill=orange!8,
        thick,
        align=center,
        minimum width=3.35cm,
        minimum height=0.95cm
    },
    arrow/.style={-{Latex[length=2.2mm]}, thick}
]

% top row
\node[box] (data) {Graph dataset\\$(A,X,Y)$};
\node[keybox, right=of data] (full) {Run full\\\textsc{WG-SRC} scaffold};
\node[keybox, right=of full] (atlas) {Compute operational\\fingerprint $\mathbf{m}(D)$};

% bottom row
\node[designbox, below=of atlas] (diagnosis) {Dataset diagnosis\\from $\mathbf{m}(D)$};
\node[designbox, left=of diagnosis] (design) {Mechanism hypotheses\\and diagnostic guidance};
\node[box, left=of design] (validate) {Follow-up model analysis\\or targeted modification};
\node[box, below=of validate] (out) {Diagnostic evidence\\for model behavior};

% arrows
\draw[arrow] (data) -- (full);
\draw[arrow] (full) -- (atlas);
\draw[arrow] (atlas) -- (diagnosis);
\draw[arrow] (diagnosis) -- (design);
\draw[arrow] (design) -- (validate);
\draw[arrow] (validate) -- (out);

% small annotations
\node[align=center, font=\scriptsize, below=0.18cm of full] {generalist probe};
\node[align=left, font=\scriptsize, below right=0.08cm and 0.18cm of atlas]
{mechanism measurement};

\end{tikzpicture}
\caption{
Mechanism atlas as a dataset diagnostic probe. \SRC{} first performs prediction
with the full scaffold and then exposes an operational fingerprint
$\mathbf{m}(D)$. The atlas is not claimed to automate model design; rather, it
produces mechanism hypotheses and post-evaluation diagnostic guidance whose support is
evaluated through aligned interventions, error shifts, and PCA--Ridge decision phase portraits on the same
measured dataset.
}
\label{fig:atlas-guided-diagnosis}
\end{figure}

\begin{table}[t]
\centering
\footnotesize
\caption{
From dataset fingerprint to post-evaluation same-dataset diagnostic guidance.
The atlas signatures generate mechanism hypotheses and diagnostic guidance for
later model analysis and dataset-specific modification on the measured dataset. The table is
a diagnostic-consistency summary, not a test-set-tuned model search: it
records whether the suggested mechanism direction agrees with aligned
interventions.
}
\label{tab:atlas-diagnostic-consistency}
\resizebox{\linewidth}{!}{%
\begin{tabular}{p{0.25\linewidth}p{0.27\linewidth}p{0.42\linewidth}}
\toprule
Atlas signature & Dataset diagnosis & Diagnostic guidance \\
\midrule
High $L_D$, low $H_D$, and positive $\Delta H_D$
& High-pass components may act as noise
& Remove, downweight, gate, or regularize high-pass blocks \\
Near-zero $R_D$ and high $C_D$
& Raw features are insufficient; class geometry is complex
& Preserve graph-derived signals and improve class-specific subspace modeling \\
High $H_D$ and negative $\Delta H_D$
& Correct nodes rely on high-pass structure
& Preserve or adaptively gate high-pass differences rather than removing them globally \\
High $Q^{\mathrm{ridge}}_D$
& Boundary decisions correct PCA failures
& Strengthen the discriminative head or improve PCA--Ridge boundary fusion \\
High $Q^{\mathrm{hard}}_D$
& Both geometry and boundary fail on many nodes
& Investigate new signal blocks, hard-node treatment, uncertainty handling, or graph rewiring \\
High $R_D$ and low-to-moderate $C_D$
& Raw features carry strong class information
& Use raw-preserving skip design, weak propagation, or propagation gating \\
\bottomrule
\end{tabular}%
}
\end{table}
